\begin{notechapter}
\chapter{Self-Supervised Representation Learning in Practice}\label{ch:contrastive-practice}
\chaptersubtitle{SimCLR, MoCo, CLIP, and the methods that need no negatives}

\begin{chapterintro}
The InfoNCE loss of \cref{ch:contrastive-learning} became a practical way to learn visual
representations without labels around 2020, when a few design choices made contrastive
pre-training competitive with supervised pre-training: strong augmentations, a projection
head, many negatives, normalized embeddings and a small temperature. This chapter
describes the main frameworks, each through the problem it solves, with verified code for
their losses, and the geometric view of what they learn: alignment of positives and uniformity
on the sphere. It ends with the methods that dispense with negatives, and with the precise
sense in which they avoid collapse. You will learn
\begin{itemize}[nosep]
  \item SimCLR: two augmented views, a projection head and the NT-Xent loss over $2B$ views;
  \item how representations are evaluated: linear probing, fine-tuning, $k$-nearest neighbours;
  \item MoCo: a queue of negatives and a momentum encoder, and why the momentum matters;
  \item CLIP: a symmetric InfoNCE between images and texts, and zero-shot classification;
  \item the alignment--uniformity decomposition of the contrastive loss;
  \item complete and dimensional collapse, and how BYOL, SimSiam, Barlow Twins and VICReg avoid
        the first without negatives.
\end{itemize}
\prerequisites{\cref{ch:contrastive-learning}; convolutional networks and transfer learning
(\cref{ch:cnn-architectures,ch:transfer-learning}); batch normalization (\cref{ch:normalization}).}
\end{chapterintro}

\section{SimCLR}\label{contrastive-practice:sec:simclr}

SimCLR (Chen, Kornblith, Norouzi and Hinton, \emph{A Simple Framework for Contrastive Learning
of Visual Representations}, 2020) has four components (\cref{contrastive-practice:fig:simclr}).
\begin{itemize}
  \item \textbf{Augmentations.} Each image $x$ of a mini-batch is transformed twice by random
        transformations $t, t'$ drawn from the same family: random crop and resize, horizontal flip,
        colour distortion, Gaussian blur. The two views form a positive pair. The authors find the
        composition of random cropping and colour distortion crucial: with crops alone, the colour
        histogram of an image identifies both of its crops, and the network can solve the task
        without learning anything about shapes.
  \item \textbf{Encoder.} A network $f$ (a ResNet-50 in the paper) maps each view to a representation
        $h = f(\tilde x)$, the vector used downstream.
  \item \textbf{Projection head.} A small MLP $g$ maps $h$ to $z = g(h)$, of dimension $128$, on which the loss is computed.
  \item \textbf{Loss.} The NT-Xent loss (normalized temperature-scaled cross-entropy), an InfoNCE over the
        $2B$ views of the batch.
\end{itemize}

\begin{figure}[ht]
\centering
\begin{tikzpicture}[
  box/.style={draw, rounded corners=2pt, minimum height=7mm, inner sep=3pt, font=\small, align=center},
  ar/.style={-{Stealth[length=1.8mm]}, thick},
  lab/.style={font=\scriptsize, inner sep=1pt}]
  \node[box, fill=gray!10] (x) at (0,0) {image $x$};
  \node[box, fill=bookaccent!12] (x1) at (2.4,0.9) {view $\tilde x_1$};
  \node[box, fill=bookaccent!12] (x2) at (2.4,-0.9) {view $\tilde x_2$};
  \draw[ar] (x) -- node[lab, above left]{$t\sim\mathcal T$} (x1); \draw[ar] (x) -- node[lab, below left]{$t'\sim\mathcal T$} (x2);
  \node[box, fill=bookexample!15] (f1) at (4.7,0.9) {encoder $f$}; \node[box, fill=bookexample!15] (f2) at (4.7,-0.9) {encoder $f$};
  \node[box, fill=bookkey!12] (h1) at (6.6,0.9) {$h_1$}; \node[box, fill=bookkey!12] (h2) at (6.6,-0.9) {$h_2$};
  \node[box, fill=bookdef!15] (g1) at (8.3,0.9) {head $g$}; \node[box, fill=bookdef!15] (g2) at (8.3,-0.9) {head $g$};
  \node[box] (z1) at (9.9,0.9) {$z_1$}; \node[box] (z2) at (9.9,-0.9) {$z_2$};
  \foreach \i in {1,2} {\draw[ar] (x\i) -- (f\i); \draw[ar] (f\i) -- (h\i); \draw[ar] (h\i) -- (g\i); \draw[ar] (g\i) -- (z\i);}
  \draw[<->, thick, bookpitfall] (z1) -- node[lab, right, text=bookpitfall]{maximize $\cos(z_1, z_2)$} (z2);
  \draw[dashed, gray] (5.6,-1.6) -- (5.6,1.6);
  \node[lab, text=gray] at (5.65,-1.85) {weights shared};
  \node[lab, align=center] at (6.6,1.75) {used\\[-2pt]downstream};
\end{tikzpicture}
\caption{SimCLR. Two random augmentations of the same image go through the same encoder and
projection head; the loss compares the projections, and the representation $h$ before the head is kept.}
\label{contrastive-practice:fig:simclr}
\end{figure}

\paragraph{The NT-Xent loss.} Number the $2B$ views so that $2k - 1$ and $2k$ come from image $k$, let $s_{ij} = \cos(z_i, z_j)$, and let $\pi(i)$ be
the partner of view $i$. Each view is an anchor whose candidates are the $2B - 1$ other views:
\[
  \ell_i = -\log\frac{e^{s_{i\pi(i)}/\tau}}{\sum_{k\neq i}e^{s_{ik}/\tau}},\qquad\mathcal L = \frac1{2B}\sum_{i=1}^{2B}\ell_i .
\]
It is the symmetric in-batch InfoNCE of \cref{contrastive-learning:sec:batch} with $N = 2B - 1$ candidates, $2B - 2$ of them
negatives. A batch is a $2B\times2B$ similarity matrix: the diagonal is excluded, each row is a softmax
classification, and the target column of row $i$ is $\pi(i)$. In code:
\begin{lstlisting}[language=Python]
import torch
import torch.nn.functional as F

def nt_xent(z, tau):
    """z: (2B, d); rows 2k and 2k+1 are the two views of image k."""
    z = F.normalize(z, dim=1)
    sim = z @ z.T / tau                       # cosines / tau
    sim.fill_diagonal_(float("-inf"))         # a view is not its own candidate
    partner = torch.arange(z.shape[0]) ^ 1    # xor: 0<->1, 2<->3, ...
    return F.cross_entropy(sim, partner)      # mean of the l_i
\end{lstlisting}
Setting the diagonal to $-\infty$ gives it probability $0$ in the softmax, which removes the term $k = i$ from the
denominator without changing the tensor's shape.

\begin{example}[Four views in the plane]\label{contrastive-practice:ex:ntxent}
Two images, $B = 2$, with unit projections at angles $0^\circ$ and $20^\circ$ (image $1$) and $100^\circ$ and $140^\circ$ (image $2$). The
cosine similarities are
\[
  S = \begin{bmatrix}1 & 0.9397 & -0.1736 & -0.7660\\ 0.9397 & 1 & 0.1736 & -0.5000\\ -0.1736 & 0.1736 & 1 & 0.7660\\ -0.7660 & -0.5000 & 0.7660 & 1\end{bmatrix}.
\]
With $\tau = 0.5$, view $1$ has $\ell_1 = -\log\frac{e^{1.8794}}{e^{1.8794} + e^{-0.3473} + e^{-1.5321}} = 0.1318$, and similarly $\ell_2 = 0.2408$, $\ell_3 = 0.3774$, $\ell_4 = 0.1188$, so
$\mathcal L = 0.2172$. View $3$ has the largest loss: its partner is $40^\circ$ away, and view $2$ is only $80^\circ$ away. The
function \texttt{nt\_xent} returns the same value.
\end{example}

\paragraph{The projection head.} The loss is computed on $z = g(h)$, but the representation kept for
downstream tasks is $h$. Chen et al.\ report that a linear classifier on $h$ is more than $10$ percentage
points more accurate than on $z$. The usual explanation: the contrastive loss makes $z$ invariant to
the augmentations, discarding information such as colour or orientation, which may be useful
downstream; the head absorbs this loss of information, and $h$ keeps it.

\paragraph{Batch size and optimization.} With in-batch negatives, the number of negatives is
tied to the batch size, and SimCLR benefits from very large batches (up to $8192$, trained with the
LARS optimizer on many accelerators), with $\tau = 0.1$ in its best ImageNet setting. A ResNet-50
pre-trained this way reaches $69.3\%$ top-1 accuracy on ImageNet with a linear classifier, and $76.5\%$ with
a $4\times$ wider network, the accuracy of the supervised ResNet-50.

\section{Evaluating representations}\label{contrastive-practice:sec:evaluation}

A self-supervised representation is judged by what it makes easy downstream.
\begin{itemize}
  \item \textbf{Linear probing.} Freeze the encoder and train a linear classifier (logistic
        regression) on its representations of a labelled data set. This measures how linearly
        separable the classes are; it is the standard ImageNet benchmark.
  \item \textbf{$k$-nearest neighbours.} Classify a test image by a vote among its nearest training
        images in representation space; no training at all.
  \item \textbf{Fine-tuning and transfer.} Use the encoder as initialization and train it on a
        downstream task, possibly with few labels, as in \cref{ch:transfer-learning}.
\end{itemize}
The contrastive loss itself is a poor guide: a lower InfoNCE does not imply better downstream
accuracy (\cref{contrastive-learning:sec:infonce}).

\section{MoCo: a queue and a momentum encoder}\label{contrastive-practice:sec:moco}

SimCLR's negatives are limited by the batch. MoCo (He, Fan, Wu, Xie and Girshick, \emph{Momentum
Contrast for Unsupervised Visual Representation Learning}, 2020) keeps a \emph{queue} of the
encoded keys of the last $K$ batches ($K = 65\,536$ in the paper) and uses them as negatives, so the
number of negatives no longer depends on the batch size. For a query $q = f_q(\tilde x)$ and its positive key
$k_+ = f_k(\tilde x')$ (normalized vectors),
\[
  \ell = -\log\frac{e^{q^\top k_+/\tau}}{e^{q^\top k_+/\tau} + \sum_{i=1}^Ke^{q^\top k_i/\tau}},
\]
with $\tau = 0.07$. The keys in the queue were computed by the key encoder at earlier steps, and no gradient
flows into them. If the key encoder were the query encoder itself, the queue would contain
keys from many different versions of a rapidly changing network, which are not comparable
with each other or with the current query. MoCo makes the key encoder a slowly moving
average of the query encoder,
\[
  \theta_k\leftarrow m\,\theta_k + (1 - m)\,\theta_q,\qquad m = 0.999,
\]
so that consecutive keys are produced by nearly the same function; with $m = 0$ (the key encoder
copied at each step) the authors report that training fails. A training step:
\begin{lstlisting}[language=Python]
def moco_step(x_q, x_k):                       # two augmented views of the same batch
    global queue                               # (K, d) normalized keys of earlier batches
    q = F.normalize(f_q(x_q), dim=1)           # queries, (B, d)
    with torch.no_grad():                      # no gradient flows into the keys
        for p_k, p_q in zip(f_k.parameters(), f_q.parameters()):
            p_k.mul_(m).add_((1 - m) * p_q)    # momentum update of the key encoder
        k = F.normalize(f_k(x_k), dim=1)       # positive keys, (B, d)
    l_pos = (q * k).sum(1, keepdim=True)       # (B, 1)
    l_neg = q @ queue.T                        # (B, K) similarities to the queue
    logits = torch.cat([l_pos, l_neg], 1) / tau
    labels = torch.zeros(len(q), dtype=torch.long)   # the positive is column 0
    loss = F.cross_entropy(logits, labels)
    opt.zero_grad(); loss.backward(); opt.step()
    queue = torch.cat([k, queue])[:K]          # enqueue the new keys, drop the oldest
    return loss
\end{lstlisting}
MoCo~v2 added SimCLR's projection head and stronger augmentations, and matched its accuracy
with batches of $256$.

\section{CLIP: contrasting images with text}\label{contrastive-practice:sec:clip}

CLIP (Radford et al., \emph{Learning Transferable Visual Models from Natural Language
Supervision}, 2021) applies the same loss across modalities: an image encoder $f_I$ and a text
encoder $f_T$ are trained on $400$ million image--caption pairs collected from the web. For a batch of $B$ pairs
$(I_i, T_i)$, with normalized embeddings, the logits are $\ell_{ij} = f_I(I_i)^\top f_T(T_j)/\tau$, and the loss is the average of two
cross-entropies over the $B\times B$ matrix, one over each row (which text matches image $i$?) and one over each
column (which image matches text $j$?):
\[
  \mathcal L = \frac12\Bigl(-\frac1B\sum_i\log\frac{e^{\ell_{ii}}}{\sum_je^{\ell_{ij}}} - \frac1B\sum_j\log\frac{e^{\ell_{jj}}}{\sum_ie^{\ell_{ij}}}\Bigr).
\]
The batch size is $32\,768$, and the temperature is learned, as a log-parameterized scale initialized to
$1/0.07$ and clipped at $100$. Because the text encoder can embed any phrase, CLIP classifies images
\emph{zero-shot}: embed the prompts ``a photo of a \{label\}'' for every class name, and assign the image to
the class whose prompt embedding has the highest cosine with the image embedding.

\section{Alignment and uniformity}\label{contrastive-practice:sec:alignment}

Wang and Isola (\emph{Understanding Contrastive Representation Learning through Alignment
and Uniformity on the Hypersphere}, 2020) describe what the contrastive loss does to normalized
representations $f(x)\in S^{d-1}$ with two quantities:
\[
  \mathcal L_{\mathrm{align}} = \E_{(x,y)\sim p_{\mathrm{pos}}}\norm{f(x) - f(y)}^2,\qquad
  \mathcal L_{\mathrm{uniform}} = \log\E_{x,x'\sim p_{\mathrm{data}}}e^{-2\norm{f(x) - f(x')}^2} .
\]
Alignment is small when positives are mapped close together; uniformity is small when the
representations spread over the sphere (the expectation is over independent pairs).

\begin{theorem}[Asymptotic decomposition; Wang and Isola]\label{contrastive-practice:thm:alignment}
For normalized representations, critic $f(x)^\top f(y)/\tau$ and $M$ negatives,
\[
  \lim_{M\to\infty}\bigl(\mathcal L_{M+1} - \log M\bigr) = -\frac1\tau\E_{p_{\mathrm{pos}}}\bigl[f(x)^\top f(y)\bigr] + \E_x\Bigl[\log\E_{x^-}\bigl[e^{f(x)^\top f(x^-)/\tau}\bigr]\Bigr].
\]
\end{theorem}

The first term equals $\frac1{2\tau}\mathcal L_{\mathrm{align}} - \frac1\tau$, since $\norm{a - b}^2 = 2 - 2a^\top b$ for unit vectors
(\cref{contrastive-practice:exr:2}); the second is minimized, under mild conditions, by representations
uniformly distributed on the sphere, which $\mathcal L_{\mathrm{uniform}}$ measures in a related form. The contrastive loss thus
trades off two goals: make positives coincide, and spread everything else out. Wang and Isola
find that directly optimizing $\mathcal L_{\mathrm{align}} + \lambda\mathcal L_{\mathrm{uniform}}$ gives representations as good as InfoNCE, and that
both quantities predict downstream accuracy.

\section{Collapse, and methods without negatives}\label{contrastive-practice:sec:collapse}

\paragraph{Two kinds of collapse.} \emph{Complete collapse} maps every input to the same vector: the
positive pairs are perfectly aligned and the representation is useless. \emph{Dimensional
collapse} is subtler: the representations vary, but only in a subspace of dimension much lower
than $d$. It is measured by the spectrum of the covariance matrix of the representations,
$\Sigma = \E\bigl[(f(x) - \mu)(f(x) - \mu)^\top\bigr]$ with $\mu = \E f(x)$: many near-zero eigenvalues indicate dimensional collapse. Negatives
prevent complete collapse, since the uniformity term penalizes it, but Jing et al.\
(\emph{Understanding Dimensional Collapse in Contrastive Self-Supervised Learning}, 2022)
observe dimensional collapse in contrastive methods too, and relate it to strong
augmentations and to the implicit regularization of training.

\paragraph{Methods without negatives.} A loss on positive pairs alone, $\norm{z_1 - z_2}^2$, is minimized by complete
collapse. Several methods nonetheless learn good representations from positives only.
\begin{itemize}
  \item \textbf{BYOL} (Grill et al., \emph{Bootstrap Your Own Latent}, 2020): an \emph{online} network with an extra
        predictor MLP $q$ predicts the projection that a \emph{target} network gives to the other view; the
        target is a momentum average of the online network, as in MoCo, and receives no gradient.
        The loss is $\norm{\bar q(z_1) - \bar z'_2}^2$ for normalized vectors, symmetrized.
  \item \textbf{SimSiam} (Chen and He, \emph{Exploring Simple Siamese Representation Learning}, 2021):
        no momentum encoder; the same network on both views, a predictor on one side and a
        \emph{stop-gradient} on the other. The authors show that removing the stop-gradient makes
        training collapse immediately.
  \item \textbf{Barlow Twins} (Zbontar et al., 2021): standardize each dimension of the two batches of
        embeddings, compute their $d\times d$ cross-correlation matrix $C$, and minimize
        $\sum_i(1 - C_{ii})^2 + \lambda\sum_{i\neq j}C_{ij}^2$: corresponding dimensions of the two views must agree, different
        dimensions must be decorrelated, which rules out complete and dimensional collapse.
  \item \textbf{VICReg} (Bardes, Ponce and LeCun, 2022): an invariance term (mean squared distance
        between the two views), a variance term (a hinge keeping the standard deviation of every
        dimension above a threshold) and a covariance term (penalizing off-diagonal
        covariances).
\end{itemize}
The first two avoid collapse through the dynamics of training, not through the loss, whose
global minimum is still the constant solution; why the asymmetric architecture and the
stop-gradient prevent the optimizer from reaching it is the subject of ongoing analysis. The last
two put explicit anti-collapse terms in the loss, playing the role that negatives play in
InfoNCE.

\begin{pitfall}[``Without negatives, representations collapse'']
The statement is true of the loss $\norm{z_1 - z_2}^2$ trained naively, and false as a general claim: BYOL and
SimSiam train without negatives and do not collapse in practice, and Barlow Twins and VICReg
prevent collapse with variance and decorrelation terms instead of negatives. Conversely,
having negatives does not rule out dimensional collapse.
\end{pitfall}

\begin{chaptersummary}
\begin{itemize}
  \item SimCLR: two augmentations per image (cropping and colour distortion are essential), a shared
        encoder, a projection head, and the NT-Xent loss over the $2B - 1$ other views; the representation before the
        head is kept.
  \item Representations are evaluated by linear probing, $k$-NN classification and fine-tuning.
  \item MoCo decouples the number of negatives from the batch size with a queue, and keeps the
        queued keys consistent with a momentum encoder ($m = 0.999$).
  \item CLIP applies a symmetric InfoNCE to image--text pairs with a learned temperature, which
        enables zero-shot classification by text prompts.
  \item With many negatives the contrastive loss decomposes into an alignment term and a
        uniformity term on the sphere.
  \item Complete collapse is prevented by negatives or by explicit variance/decorrelation terms
        (Barlow Twins, VICReg); BYOL and SimSiam avoid it through asymmetry and stop-gradients.
        Dimensional collapse can occur in all of them.
\end{itemize}
\end{chaptersummary}

\section*{Exercises}
\addcontentsline{toc}{section}{Exercises}

\begin{exercise}\label{contrastive-practice:exr:1}
In \cref{contrastive-practice:ex:ntxent}, compute $\ell_3$. How does $\mathcal L$ change if $\tau\to0$? If $\tau\to\infty$?
\solutionlink{contrastive-practice:exr:1}
\end{exercise}

\begin{exercise}\label{contrastive-practice:exr:2}
Show that for unit vectors $\norm{a - b}^2 = 2 - 2a^\top b$, and deduce that
\[
  -\frac1\tau\E[f(x)^\top f(y)] = \frac1{2\tau}\mathcal L_{\mathrm{align}} - \frac1\tau .
\]
\solutionlink{contrastive-practice:exr:2}
\end{exercise}

\begin{exercise}\label{contrastive-practice:exr:3}
How many negatives does each anchor have in SimCLR with $B = 4096$? In MoCo with $B = 256$ and $K = 65\,536$? How
many gradient-carrying embeddings does each method compute per step?
\solutionlink{contrastive-practice:exr:3}
\end{exercise}

\begin{exercise}\label{contrastive-practice:exr:4}
With the momentum update $\theta_k\leftarrow m\theta_k + (1 - m)\theta_q$ and $m = 0.999$, the key encoder is an exponential moving average of the
query encoder. Give the weight of the query parameters of $1\,000$ steps ago, and the number of steps over
which the weights sum to about $63\%$.
\solutionlink{contrastive-practice:exr:4}
\end{exercise}

\begin{exercise}\label{contrastive-practice:exr:5}
Compute the Barlow Twins cross-correlation matrix $C$ and the loss when every embedding is the same vector
(complete collapse), assuming the standardization divides by the standard deviation plus a small
$\varepsilon > 0$. What is the loss for two views whose standardized embeddings are equal and whose dimensions
are uncorrelated?
\solutionlink{contrastive-practice:exr:5}
\end{exercise}
\end{notechapter}
